\section{6. Differentialgleichungen}

\subsection{Teil 1}

\Def In \key{Differentialgleichung} treten gesuchte Funktion sowie Ableitungen auf,
z.B $y'(x) = 69x \cdot y(x) + 420$. $x$ ist \key{unabhängige Variable} und $y$ heisst
\key{abhängige Variable}

\Def Eine DGL der Form
$a_n(x)y^{(n)} + \ldots + a_0(x)y(x) = s(x)$
heisst \key{lineare} DGL mit \key{Koeffizienten} $a_i(x)$. Die \key{Ordnung} der DGL ist
die maximale Ordnung der Ableitungen. $s(x)$ heisst \key{Störfunktion}. Falls $s(x) = 0$
\key{homogen}, sonst \key{inhomogen}.
Falls alle $a_i(x)$ konstant, DGL \key{mit konstanten Koeffizienten}

\Satz \key{Superpositionsprinzip}: $y_1(x)$ und $y_2(x)$ Lösungen einer linearen, homogenen DGL.
$y(x) = C_1 y_1(x) + C_2 y_2(x)$, $C_1, C_2 \in \RR$ ist auch eine Lösung.

\Satz \key{Fundamentallösungen}: Lineare, homogene DGL der Ordnung $n$ hat $n$
\key{Fundamentallösungen} $y_1(x), \ldots, y_n(x)$ so dass die allgemeine Lösung gegeben ist durch
$y(x) = C_1 y_1(x) + \ldots + C_n y_n(x)$, $C_1, \ldots, C_n \in \RR$

\subsection{Lösen einer linearen, homogenen DGL}

Gegeben $a_n u^{(n)}(x) + \ldots + a_0 u(x) = 0$. Der \key{Ansatz} $u(x) = e^{\lambda x}$
(Eulerscher Ansatz) führt uns zur \key{charakteristischen Gleichung}
$a_n \lambda^n + \ldots + a_1 \lambda + a_0 = 0$ und das \key{charakteristische Polynom}
$p(\lambda) = a_n \lambda^n + \ldots + a_1 \lambda + a_0$.

Nun für die allgemeine Lösung der DGL:
\vspace{-0.6em}
\begin{itemize}
    \item Falls alle $r$ Nullstellen $\lambda_i$ von $p(\lambda)$ reell mit Multiplizität $m_i$:
    $u(x) = \sum_{i=1}^{r} \left(\sum_{p=0}^{m_i-1} C_{ip}x^pe^{\lambda_i x}\right)$
    \item Falls $r$ reelle Nullstelle und $s$ \textit{Paare} komplexer Nullstellen ($a_j \pm ib_j$)
    mit Multiplizität $m_i$:
    $u(x) = \sum_{i=1}^{r}\left(\sum_{p=0}^{m_i - 1}C_{ip} x^p e^{\lambda_i x}\right) + \sum_{j=1}^{s} \left(\sum_{q=0}^{m_j-1}(A_{jq} x^q e^{a_j x} \sin(b_j x) + B_{jq} x^q e^{a_j x} \cos(b_j x))\right)$
\end{itemize}

In der allgemeinen Lösung einer DGL $n$-ter Ordnung tauchen \key{$n$ Konstanten} auf.

Wollen wir also aus der allgemeinen Lösung \textit{eine} Funktion auswählen, so müssen wir \key{$n$ Informationen vorgeben}

$\to$ Zwei Strategien
\vspace{-0.6em}
\begin{itemize}
    \item \key{Randwertproblem}: Wir geben die Informationen an \textit{zwei unterschiedlichen Stellen} vor, z.B., im fall $n = 2$,
    $$u(t_0) = u_0 \hspace{20pt} u(t_1) = u_1 \hspace{20pt} t_0 \ne t_1$$
    \item \key{Anfangswertproblem}: Wir geben an \textit{einer Stelle} die Informationen vor, z.B., im fall $n = 2$,
    $$u(t_0) = u_0 \hspace{20pt} u'(t_0) = v_0$$
\end{itemize}

\subsection{Aufstellen von Differentialgleichungen}

\subsubsection{Typische Situation}

Gegeben Geburtenrate $B(t)$ und Sterberate $T(t)$, wie entwickelt sich Bevölkerung $y(t)$?

Wie verändert sich $y(t)$ in einem kleinen Zeitintervall $\Delta t$?
\vspace{-0.6em}
\begin{itemize}
    \item Zuwachs: $B(t) \cdot \Delta t$
    \item Abnahme: $T(t) \cdot \Delta t$
    \item Nettobilanz: $y(t + \Delta t) = y(t) + B(t) \cdot \Delta t - T(t) \cdot \Delta t$
\end{itemize}

Was ist die Differenz zwischen nahe den Werten, falls die Zeiten nahe beieinander sind?
$$\Delta y = y(t + \Delta t) - y(t) = B(t) \cdot \Delta t - T(t) \cdot \Delta t $$

Was ist die proportionale Veränderung?
$$\frac{\Delta y}{\Delta t} = \frac{y(t + \Delta t) - y(t)}{\Delta t} = B(t) - T(t)$$

Was passiert im Grenzwert $\Delta t \to 0$?
$$\frac{dy}{dt} = \lim_{\Delta t \to 0} \frac{\Delta y}{\Delta t} = \lim_{\Delta t \to 0}\frac{y(t + \Delta t) - y(t)}{\Delta t} = B(t) - T(t)$$

Also ist die gesuchte Differentialgleichung
$$y'(t) = B(t) - T(t)$$

\subsubsection{Mechanische Probleme}

Hier ist zu beachten, dass das Problem meist über die involvierten Kräfte beschieben werden kann.
Für jede Kraft $F$ gilt
$$F = m \cdot y'' = m \cdot a$$
wobei $y$ die Position des betrachteten Körpers und $a$ dessen Beschleunigung sind.

\subsection{Teil 2}

\begin{itemize}
    \item Trennung der Variablen
    \item Substitution
    \item Variation der Konstanten (lineare, inhomogene DGL)
\end{itemize}
\textit{Lineare, inhomogene} DGL \textit{beliebiger} Ordnung:
\vspace{-0.6em}
\begin{itemize}
    \item Geeignete Ansätze für die partikuläre Lösung
\end{itemize}

\Satz \key{Trennung der Variablen}:
$$y'(x) = f(x) \cdot g(y) \implies \int \frac{1}{g(y)} dy = \int f(x) dx$$

\Satz \key{Lineare Substitution} Falls DGL der Form
$y'(x) = f(ax + by(x) + c)$, mit $a, b, c \in \RR$, dann erhalten wir durch die
Substitution $u = ax + by(x) + c$ eine DGL mit trennbaren Variablen

\Def Eine DGL, wo die unab. var. $x$ nie vorkommt, heisst \key{autonom}

\Def Eine DGL folgender Form heisst \key{homogen in den Variablen}:
$$y' = g\left(\frac{y}{x}\right)$$

\Satz \key{Substitution bei in den Variablen homogenen DGL} Falls DGL homogen in den Variablen
erhalten wir durch die Substitution $u = \frac{y}{x}$ eine DGL mit trennbaren Variablen

\Def Falls DGL der Form
$$a_n(x)y^{(n)}(x) + \cdot + a_0(x)y(x) = s(x) \ne 0$$
heisst $a_n(x)y^{(n)}(x) + \cdot + a_0(x)y(x) = 0$ die \key{dazugehörige homogene DGL}

\Satz Ist $y_p(x)$ eine \key{partikuläre Lösung} einer \textit{linearen, inhomogenen} DGL und
$y_h(x)$ die allgemeine Lösung der dazugehörige homogene DGL, so ist die allgemeine Lösung der
\textit{linearen, inhomogenen} DGL $y(x) = y_h(x) + y_p(x)$

\Satz \key{Variation der Konstanten}: Gegeben eine \textit{inhomogene, lineare} DGL erster Ordnung.
Dann erhalten wir einen \textit{Ansatz} für die partikuläre Lösung, indem wir der allgemeinen
Lösung der dazugehörigen homogenen DGL die auftretende Konstante durch
\textit{eine Funktion in der unab. Var. ersetzen}

\subsubsection{Ansätze für die partikuläre Lösung}

\begin{minipage}{\linewidth}
\begin{tabular}{p{0.47\linewidth} p{0.47\linewidth}}
\textbf{Störfunktion $s(t)$} & \textbf{Ansatz $y_p(t)$} \\
\hline
\key{Polynom}: $a_0 + a_1 t + \cdots + a_n t^n$ & $C_0 + C_1 t + \cdots + C_n t^n$ \\
Spezialfall: $0$ ist $m$-fache NS d. char. Pol. & $(C_0 + \cdots + C_n t^n)\,t^m$ \\
\hline
\key{Exponentialfkt.}: $A e^{kt}$ & $C e^{kt}$ \\
Spezialfall: $k$ ist $m$-fache NS d. char. Pol. & $(C e^{kt})\,t^m$ \\
\hline
\key{Schwingung}: \\ $A \sin(\omega t) + B \cos(\omega t)$ & $C_1 \sin(\omega t) + C_2 \cos(\omega t)$ \\
Spezialfall: $i\omega$ ist $m$-fache NS d. char. Pol. & $(C_1 \sin(\omega t) + C_2 \cos(\omega t))\,t^m$
\end{tabular}
\end{minipage}
